\section{初唐的共同语言与不均等的经验}
\label{sec:early-tang-society-institutions}

贞观与永徽时期的律令、官制和都城常被作为一个“成熟帝国”的证据。更接近材料的读法，是把它们看作使不同人能够同官府打交道的共同语言：户婚、刑名、官司与仓库有了可援引的分类，官员也能以同一套文书解释职责。可援引不等于可抵达。唐律疏议保存的是规范及其官方解释，敦煌、吐鲁番文书保存的是特定地点的实践；二者之间的距离，正是地方家庭、书吏和官员变通制度的空间。%
\footnote{贞观律令与《唐律疏议》见 ST-627-ZHENGUAN-CODE、ST-653-TANGLU-SHUYI，核心材料为《旧唐书·刑法志》《唐律疏议》《通典》及令典残卷。不能从成文法直接推出全国司法实践。}

长安也不是只由皇帝和宰相居住的舞台。635年阿罗本获准译经、664年玄奘死后译场继续活动，说明僧侣、译者、施主、官府和道路把不同语言的宗教知识接到都城；它们并不证明国家为所有社群提供同等空间。景教碑、僧传和寺院题记都是各自的自我表述，适于显示网络与仪式，不能据此计算信徒人数，或把一座都城的情况外推为整个乡村社会。%
\footnote{阿罗本与玄奘的节点见 ST-635-NESTORIAN、ST-664-XUANZANG-DEATH，依据《旧唐书》相关记载、景教碑、《大慈恩寺三藏法师传》及《续高僧传》。文本成书、碑刻立置和实际活动的日期需要分别辨认。}

皇室继承同样穿过这些制度，而非停留在制度之外。李承乾被废、李治继位，后来武后进入皇后与临朝位置，表明法典、礼仪和官僚网络可以提供争论的格式，却不能预先消除继承中的暴力与排除。初唐国家最值得追问的不是它是否已经“强大”，而是它如何把法律、城市宗教、家族地位和边防军功暂时接在一起；这个接合让帝国能够行动，也使被排除者和地方差异继续存在于文本之外。%
\footnote{承乾案、高宗即位及武后立后见 ST-643-CHENGQIAN、ST-649-TAIZONG-DEATH、ST-655-WU-EMPRESS，依据两《唐书》本纪、列传及《资治通鉴》。宫廷案件由后来的胜者保存，动机和言辞不能作确定复原。}
